\chapter{Static Routing, IPv4 and IPv6}\label{ch:static}
\bp{3.2}

\begin{outcomes}
By the end of this chapter you will be able to:
\begin{tight}
  \item \textbf{Configure} the four kinds of static route the blueprint names:
        default, network, host and floating.
  \item \textbf{Choose} correctly between a next-hop address, an exit interface,
        and both --- a choice that matters more than it looks.
  \item \textbf{Configure} static routing for IPv6 alongside IPv4.
  \item \textbf{Troubleshoot} a static route that is absent from the table, or
        present and not working.
\end{tight}
\end{outcomes}

\begin{prereq}
\chref{routingtable} --- longest prefix match and administrative distance are
assumed throughout. \chref{ipv6} for the v6 half.
\end{prereq}

\begin{keyterms}
static route $\bullet$ default route $\bullet$ network route $\bullet$ host route
$\bullet$ floating static $\bullet$ next hop $\bullet$ exit interface $\bullet$
fully specified route $\bullet$ recursive lookup $\bullet$ gateway of last resort
\end{keyterms}

\begin{examnote}[The verb is \emph{troubleshoot}, and the sub-topics are a checklist]
Topic \tp{3.2} is ``troubleshoot IPv4 and IPv6 static routing'' and then lists
exactly four kinds: default route, network route, host route, floating static.
Expect to be shown a table and a symptom rather than asked to recite syntax ---
but you cannot diagnose what you cannot write, so learn the syntax first.
\end{examnote}

\section{The four kinds}

\begin{longtable}{@{}L{2.8cm}L{4.4cm}L{7.8cm}@{}}
\toprule
\rowcolor{primary!15}
\textbf{Kind} & \textbf{Prefix} & \textbf{What it is for} \\
\midrule
\endhead
\textbf{Network} & A real subnet, e.g.\ \cmd{/24} &
The ordinary case: reach that network via this neighbour \\
\rowcolor{lightbg}
\textbf{Default} & \cmd{0.0.0.0/0} &
``Everything I do not otherwise know.'' Points at the internet or the core \\
\textbf{Host} & \cmd{/32} (or \cmd{/128}) &
One specific address. Used to steer a single server down a particular path \\
\rowcolor{lightbg}
\textbf{Floating} & Any of the above, with a raised distance &
A backup that stays out of the table until the preferred route disappears \\
\bottomrule
\end{longtable}

\begin{config}{all four, on IPv4}
! Network route -- reach 10.9.9.0/24 through the neighbour at 10.255.0.2.
ip route 10.9.9.0 255.255.255.0 10.255.0.2
!
! Default route -- the gateway of last resort.
ip route 0.0.0.0 0.0.0.0 203.0.113.1
!
! Host route -- this one server goes a different way.
ip route 10.9.9.50 255.255.255.255 10.255.1.2
!
! Floating static -- distance 200 keeps it out of the table while
! OSPF (distance 110) has a route. It appears only if OSPF loses one.
ip route 10.1.0.0 255.255.0.0 10.255.2.2 200
\end{config}

\begin{conceptbox}[Why the host route wins]
All four routes above can match \cmd{10.9.9.50}. The \cmd{/32} host route is the
longest prefix, so it is used --- ahead of the \cmd{/24}, ahead of the default.
Longest prefix match decides before distance is even considered
(\chref{routingtable}).
\end{conceptbox}

\section{Next hop, exit interface, or both}

This is the choice that produces the most subtle static-routing faults.

\begin{longtable}{@{}L{4.4cm}L{5.0cm}L{5.6cm}@{}}
\toprule
\rowcolor{primary!15}
\textbf{Form} & \textbf{Example} & \textbf{Behaviour} \\
\midrule
\endhead
\textbf{Next hop only} & \cmd{ip route 10.9.9.0 255.255.255.0 10.255.0.2} &
The router must look up how to reach \cmd{10.255.0.2} --- a \term{recursive
lookup}. Correct on any link \\
\rowcolor{lightbg}
\textbf{Exit interface only} & \cmd{ip route 10.9.9.0 255.255.255.0 Serial0/0/0} &
No recursion. Fine on point-to-point links; \textbf{dangerous on Ethernet} \\
\textbf{Fully specified} & \cmd{ip route 10.9.9.0 255.255.255.0 Gi0/1 10.255.0.2} &
Both. No recursion and no ambiguity. The safest form on Ethernet \\
\bottomrule
\end{longtable}

\begin{alertbox}[Exit-interface-only routes on Ethernet are a trap]
Ethernet is multi-access: many devices share the segment. A route that names only
an exit interface tells the router ``send it out here'' without saying to whom, so
the router must ARP for \emph{every destination address} it forwards, hoping
something answers. On a link where proxy ARP is enabled it appears to work; where
it is not, it silently fails; and either way the ARP table fills with entries that
should not exist.

On a serial point-to-point link there is only one possible recipient, so the form
is harmless. On Ethernet, always give the next hop --- alone or fully specified.
\end{alertbox}

\begin{verify}{R1}{show ip route static}
      10.0.0.0/8 is variably subnetted, 5 subnets, 3 masks
S        10.1.0.0/16 [200/0] via 10.255.2.2
S        10.9.9.0/24 [1/0] via 10.255.0.2
S        10.9.9.50/32 [1/0] via 10.255.1.2
S*    0.0.0.0/0 [1/0] via 203.0.113.1
\end{verify}

Note the floating static at \cmd{[200/0]} \emph{is} shown here --- meaning OSPF
currently has no route for \cmd{10.1.0.0/16}, so the backup has taken over. When
OSPF is healthy that line is absent from the table entirely.

\section{Static routing for IPv6}

The syntax mirrors IPv4, with one wrinkle worth learning.

\begin{config}{IPv6 static routes}
ipv6 unicast-routing
!
! Network route via a global next-hop address.
ipv6 route 2001:db8:acad:20::/64 2001:db8:acad:99::2
!
! Default route.
ipv6 route ::/0 2001:db8:acad:99::2
!
! Host route.
ipv6 route 2001:db8:acad:20::50/128 2001:db8:acad:99::2
!
! Floating static.
ipv6 route 2001:db8:acad:30::/64 2001:db8:acad:98::2 200
!
! Via a LINK-LOCAL next hop -- the exit interface is mandatory,
! because fe80:: addresses are only unique per link and the router
! has no way to know which interface you mean.
ipv6 route 2001:db8:acad:40::/64 GigabitEthernet0/1 fe80::2
\end{config}

\begin{conceptbox}[Why link-local next hops need the interface]
A link-local address is valid on one link only, and the same \cmd{fe80::2} may
legitimately exist on three different interfaces of the same router. Naming the
interface removes the ambiguity. IOS rejects a link-local next hop without one.

This matters because routing protocols, including OSPFv3, use link-local
addresses as their next hop --- so link-local next hops are the normal case in
IPv6, not an oddity.
\end{conceptbox}

\begin{verify}{R1}{show ipv6 route static}
IPv6 Routing Table - default - 7 entries
Codes: C - Connected, L - Local, S - Static, U - Per-user Static route
S   ::/0 [1/0]
     via 2001:DB8:ACAD:99::2
S   2001:DB8:ACAD:20::/64 [1/0]
     via 2001:DB8:ACAD:99::2
S   2001:DB8:ACAD:40::/64 [1/0]
     via FE80::2, GigabitEthernet0/1
\end{verify}

\section{Floating statics, properly}

\begin{worked}[A backup path that only activates on failure]
\cmd{R1} reaches \cmd{10.1.0.0/16} through OSPF over its primary link, and has a
slow secondary link to the same destination.

\begin{tightnum}
  \item OSPF's route has distance 110.
  \item Configure the backup with a distance \emph{above} 110:
        \cmd{ip route 10.1.0.0 255.255.0.0 10.255.2.2 200}.
  \item While OSPF is up, only the OSPF route is installed. \cmd{show ip route}
        shows \cmd{O 10.1.0.0/16 [110/20]} and no static entry at all.
  \item The primary link fails. OSPF withdraws its route. The static route, now
        the only candidate, is installed as \cmd{S 10.1.0.0/16 [200/0]}.
  \item The primary returns. OSPF's route at 110 beats 200 and displaces the
        static again.
\end{tightnum}

The number 200 is not magic --- it just has to exceed the distance of whatever it
is backing up. Back up a \emph{static} route (distance 1) and any value from 2
upwards will do.
\end{worked}

\begin{pitfall}
\begin{tight}
  \item \textbf{A floating static with too low a distance.} Set it below the
        dynamic protocol's and it takes over permanently, which is the opposite
        of a backup.
  \item \textbf{Exit-interface-only routes on Ethernet.} See the warning above.
  \item \textbf{Forgetting the reverse route.} A static route is one-directional.
        Both ends need one, and forgetting this is the commonest static-routing
        fault of all.
  \item \textbf{A next hop that is not directly reachable.} The route will not be
        installed at all; the router silently discards it. \cmd{show ip route}
        simply has no entry, which looks like you never typed it.
  \item \textbf{Omitting the interface on an IPv6 link-local next hop.} Rejected
        outright.
\end{tight}
\end{pitfall}

\begin{breakfix}[The static route was typed and it is not in the table.]
An engineer configures \cmd{ip route 10.9.9.0 255.255.255.0 10.255.0.2} and
\cmd{show ip route static} shows nothing. No error was displayed when the command
was entered.

\begin{verify}{R1}{show ip route 10.255.0.2}
% Network not in table
\end{verify}

\textbf{Diagnosis.} A static route is only installed if the router can resolve its
next hop. \cmd{10.255.0.2} is not reachable --- there is no connected subnet
covering it --- so the route is accepted into the configuration and never enters
the routing table. This is why the command produced no error: the syntax was
perfect and the topology was not.

Check the interface that should reach it:

\begin{verify}{R1}{show ip interface brief | include Gigabit}
GigabitEthernet0/1     10.255.0.1      YES manual administratively down down
\end{verify}

\textbf{Fix.} The interface is shut. \cmd{no shutdown} on \cmd{Gi0/1}, and the
connected route for \cmd{10.255.0.0/30} appears, the next hop resolves, and the
static route installs itself immediately. Confirm with \cmd{show ip route static}
--- and remember the general rule: \textbf{a static route missing from the table
is nearly always an unreachable next hop}, not a typing error.
\end{breakfix}

\begin{lab}{Four routes, one failure}{Packet Tracer}
\textbf{Topology.} Three routers --- \cmd{R1}, \cmd{R2}, \cmd{R3} --- with
\cmd{R1} connected to \cmd{R2} by a fast link and to \cmd{R3} by a slow one, and
\cmd{R2} and \cmd{R3} both connected to a LAN at the far side. A LAN on \cmd{R1}.

\begin{tasks}
  \item Address everything. Confirm connected routes on all three and that
        neighbours ping.
  \item On \cmd{R1}, add a network route to the far LAN via \cmd{R2}. Test from
        the \cmd{R1} LAN --- it will fail. Diagnose why, then add the reverse
        route and re-test.
  \item Add a default route on \cmd{R2} and \cmd{R3} pointing back at \cmd{R1}.
        Note the \cmd{S*} and the gateway of last resort line.
  \item Add a host route on \cmd{R1} for one server in the far LAN, via
        \cmd{R3}. Prove with \cmd{show ip route <address>} that it is preferred
        over the network route.
  \item Add a floating static via \cmd{R3} for the far LAN with distance 5.
        Confirm it is \emph{not} in the table.
  \item Shut \cmd{R1}'s link to \cmd{R2}. Confirm the floating static appears and
        traffic still flows. Restore the link and watch it withdraw.
  \item \textbf{Extension.} Repeat steps 2 and 5 for IPv6, including one route via
        a link-local next hop.
\end{tasks}
\end{lab}

\begin{checkpoint}
\begin{tightnum}
  \item You configure a static route and it never appears in the table, with no
        error message. What is the most likely cause and which command confirms
        it?
  \item A floating static backing up an OSPF route is configured with distance
        100. What will happen, and why is it wrong?
  \item Why must an IPv6 static route using a link-local next hop also name an
        exit interface?
\end{tightnum}
\end{checkpoint}

\begin{chaptersummary}
\begin{tight}
  \item Four kinds: network, default (\cmd{0.0.0.0/0}), host (\cmd{/32}) and
        floating. The host route wins on longest prefix match.
  \item Give a next hop on Ethernet --- alone or fully specified. Exit-interface
        only belongs on point-to-point links.
  \item IPv6 mirrors IPv4, but a link-local next hop \emph{must} name the exit
        interface.
  \item A floating static needs a distance above whatever it backs up: over 110
        for OSPF, over 1 for another static.
  \item A static route absent from the table almost always means an unresolvable
        next hop. Routes are one-way; configure both ends.
\end{tight}
\end{chaptersummary}

\begin{reviewq}
  \item \textbf{(MCQ)} Which command creates a default route via
        \cmd{203.0.113.1}?
        \begin{tight}
          \item[(a)] \cmd{ip route 0.0.0.0 255.255.255.255 203.0.113.1}
          \item[(b)] \cmd{ip route 0.0.0.0 0.0.0.0 203.0.113.1}
          \item[(c)] \cmd{ip default-gateway 203.0.113.1}
          \item[(d)] \cmd{ip route default 203.0.113.1}
        \end{tight}
  \item \textbf{(MCQ)} A floating static backs up an OSPF route. Which distance
        is suitable? (a)~0 (b)~1 (c)~100 (d)~200
  \item \textbf{(Short)} Explain why an exit-interface-only static route is
        acceptable on a serial link but risky on Ethernet.
  \item \textbf{(Short)} A host route and a network route both match a
        destination. Which is used, and does administrative distance play any
        part?
  \item \textbf{(Applied)} \cmd{R1} must reach \cmd{192.168.80.0/24} via
        \cmd{R2} normally and via \cmd{R3} if that path fails, and must send
        everything else to the internet via \cmd{203.0.113.1}. Write the complete
        configuration and state which \cmd{show} output would prove the backup is
        correctly dormant.
\end{reviewq}
